\documentclass[a4paper,12pt]{article}

\usepackage[french]{babel}
\usepackage[T1]{fontenc}
\usepackage[utf8]{inputenc}
\usepackage{amsmath, amssymb}
\usepackage{geometry}
\usepackage{tikz}
\usepackage{pgfplots}
\pgfplotsset{compat=1.18}

\geometry{top=2cm, bottom=2cm, left=2cm, right=2cm}

\begin{document}

\begin{center}\Large\textbf{Épreuve anticipée de mathématiques}\end{center}
\begin{center}\textbf{Voie générale : candidats suivant l'enseignement de spécialité de mathématiques.}\end{center}

\begin{center}
\textbf{Durée : 2 heures. L'usage de la calculatrice n'est pas autorisé.}
\end{center}

\section*{PREMIÈRE PARTIE : AUTOMATISMES -- QCM (6 pts)}

\textit{Pour cette première partie, aucune justification n'est demandée et une seule réponse est possible par question. Pour chaque question, reportez son numéro sur votre copie et indiquez votre réponse.}

\begin{enumerate}

\item Si $a > 0$ alors :
\begin{itemize}
\item[A.] $a = a^2$
\item[B.] $a \leq a^2$
\item[C.] $a \geq a^2$
\item[D.] $a \leq a^2$ ou $a^2 \geq a$ Cela dépend de la valeur de $a$.
\end{itemize}

\item Dans une classe de 30 élèves, 12 sont des filles. La proportion de filles est :
\begin{itemize}
\item[A.] $12\%$
\item[B.] $25\%$
\item[C.] $30\%$
\item[D.] $40\%$
\end{itemize}

\item Le prix d'un article est multiplié par $1,11$. Cela signifie que l'article a subi :
\begin{itemize}
\item[A.] une hausse de $1,11\%$
\item[B.] une hausse de $11\%$
\item[C.] une baisse de $1,11\%$
\item[D.] une baisse de $11\%$
\end{itemize}

\item On donne ci-dessous la courbe représentative d'une fonction $f$ définie sur $[-4;4]$.

\begin{center}
\begin{tikzpicture}[scale=0.8]
    \begin{axis}[
        axis lines = middle,
        xlabel = $x$,
        ylabel = $y$,
        xmin = -4.5, xmax = 5.5,
        ymin = -2.5, ymax = 3.5,
        xtick = {-4,-3,-2,-1,0,1,2,3,4,5},
        ytick = {-2,-1,0,1,2,3},
        grid = both,
        grid style = {gray!30},
        width = 12cm,
        height = 8cm,
        samples = 200,
        smooth
    ]
        % Courbe passant par (1;2) et par (-2;0) et (3;0)
        % f(x) = -0.2x^3 + 0.1x^2 + 1.2x + 1.2
        % Vérification : f(1) = -0.2 + 0.1 + 1.2 + 1.2 = 2.3 (trop haut)
        % Ajustement pour f(1)=2 : f(x) = -0.2x^3 + 0.2x^2 + 1.2x + 1.2
        % f(1) = -0.2 + 0.2 + 1.2 + 1.2 = 2.4 (encore trop haut)
        % Correction : f(x) = -0.2x^3 + 0.2x^2 + 1.0x + 1.0
        % f(1) = -0.2 + 0.2 + 1.0 + 1.0 = 2.0 (ok)
        % f(-2) = -0.2*(-8) + 0.2*4 + 1.0*(-2) + 1.0 = 1.6 + 0.8 - 2 + 1 = 1.4 (pas 0)
        % Pour avoir f(-2)=0 et f(3)=0 et f(1)=2, utilisons une fonction polynomiale de degré 3 passant par ces points
        % Par calcul : f(x) = -0.13333x^3 - 0.2x^2 + 1.3333x + 1.0667
        % f(1) = -0.13333 - 0.2 + 1.3333 + 1.0667 = 2.0667 (proche)
        % Version plus simple : courbe passant par (-2;0), (1;2), (3;0)
        % f(x) = -0.25(x+2)(x-3)(x+1) + 0.5
        \addplot[thick, domain=-4:5] { -0.25*(x+2)*(x-3)*(x+1)-0.5  };
        
        % Point à l'abscisse 1
        \addplot[only marks, mark=*, mark size=2] coordinates {(1,2)};
        
    \end{axis}
\end{tikzpicture}
\end{center}

$f(0)$ est égal à :
\begin{itemize}
\item[A.] $-1$
\item[B.] $1$
\item[C.] $0$
\item[D.] $2$
\end{itemize}

\item On a représenté ci-dessous les effectifs des différentes LV1 dans un établissement.

\begin{center}
\begin{tikzpicture}
    \begin{axis}[ybar, ymin=0, ymax=100, width=10cm, height=6cm,
        xtick={1,2,3,4}, xticklabels={Anglais, Espagnol, Allemand, Italien},
        ylabel={Effectifs}, xlabel={Langues}]
        \addplot coordinates {(1,85) (2,50) (3,30) (4,20)};
    \end{axis}
\end{tikzpicture}
\end{center}

La LV1 la plus suivie dans cet établissement est :
\begin{itemize}
\item[A.] l'espagnol
\item[B.] l'allemand
\item[C.] l'anglais
\item[D.] l'italien
\end{itemize}

\item On lance trois fois un dé équilibré. La probabilité d'obtenir au moins un 6 est environ $0,19$. On peut alors affirmer que la probabilité de n'obtenir aucun 6 lors des trois lancers est égale à :
\begin{itemize}
\item[A.] $0,81$
\item[B.] $0,91$
\item[C.] $1,19$
\item[D.] $0,19$
\end{itemize}

\item On interroge un groupe de 300 personnes sur le thème de la lecture. Les réponses sont consignées dans le tableau ci-dessous :

\[
\begin{array}{|c|c|c|c|}
\hline
& \text{Lit au moins une heure par semaine} & \text{Lit moins d'une heure par semaine} & \text{Total} \\
\hline
\text{Moins de 30 ans} & 20 & 80 & 100 \\
\hline
\text{Entre 30 et 50 ans} & 80 & 40 & 120 \\
\hline
\text{50 ans et plus} & 70 & 10 & 80 \\
\hline
\text{Total} & 170 & 130 & 300 \\
\hline
\end{array}
\]

On choisit une personne de ce groupe au hasard et on définit les événements suivants : $L$ « la personne lit au moins une heure par semaine », $A_1$ « la personne a moins de 30 ans », $A_2$ « la personne a entre 30 et 50 ans » et $A_3$ « la personne a plus de 50 ans ».

$\dfrac{70}{80}$ correspond à la valeur de :
\begin{itemize}
\item[A.] $P(A_3)$
\item[B.] $P(L \cap A_3)$
\item[C.] $P_{A_3}(L)$
\item[D.] $P_L(A_3)$
\end{itemize}

\item Un ingénieur calcule la vitesse maximale $V$ d'un train à grande vitesse en km/h. Il est plus vraisemblable qu'il trouve :
\begin{itemize}
\item[A.] $V = 25$
\item[B.] $V = 90$
\item[C.] $V = 350$
\item[D.] $V = 1240$
\end{itemize}

\item La fonction $f$ définie sur $\mathbb{R}$ par $f(x) = 3x - 1$ admet pour tableau de signe :
\begin{itemize}
\item[A.] 
$\begin{array}{c|ccc}
x & -\infty & \frac{1}{3} & +\infty \\ \hline
f(x) & - & 0 & +
\end{array}$
\item[B.]
$\begin{array}{c|ccc}
x & -\infty & \frac{1}{3} & +\infty \\ \hline
f(x) & + & 0 & -
\end{array}$
\item[C.]
$\begin{array}{c|ccc}
x & -\infty & 0 & +\infty \\ \hline
f(x) & - & 0 & +
\end{array}$
\item[D.]
$\begin{array}{c|ccc}
x & -\infty & 0 & +\infty \\ \hline
f(x) & + & 0 & -
\end{array}$
\end{itemize}

\item Un potager de 30 m$^2$ représente $\dfrac{2}{7}$ de la surface d'un jardin. Le jardin a une surface totale égale à :
\begin{itemize}
\item[A.] 27 m$^2$
\item[B.] 105 m$^2$
\item[C.] 210 m$^2$
\item[D.] 140 m$^2$
\end{itemize}

\item Lors d'un été très chaud, le niveau d'une nappe phréatique baisse de $30\%$ au mois de juillet puis de $20\%$ au mois d'août. Le niveau a globalement baissé de :
\begin{itemize}
\item[A.] $6\%$
\item[B.] $44\%$
\item[C.] $50\%$
\item[D.] $56\%$
\end{itemize}

\item On considère la courbe d'équation $y = \dfrac{6}{x}$. Déterminer le point qui n'appartient pas à cette courbe :
\begin{itemize}
\item[A.] $M(2;3)$
\item[B.] $N(3;3)$
\item[C.] $P(1;6)$
\item[D.] $Q\left(\dfrac{1}{2};12\right)$
\end{itemize}

\end{enumerate}

\newpage

\section*{DEUXIÈME PARTIE (14 pts)}

\subsection*{Exercice 1}

Une balle est lancée en l'air à l'instant $t = 0$ (exprimé en secondes). Sa hauteur $h(t)$ (exprimée en mètres) est donnée par :
\[
h(t) = -4t^2 + 16t + 2
\]

\begin{enumerate}
\item On donne ci-dessous la courbe représentative de la fonction $h$.

\begin{center}
\begin{tikzpicture}[scale=0.7]
    \begin{axis}[
        axis lines = middle,
        xlabel = $t$,
        ylabel = $h(t)$,
        xmin = -0.5, xmax = 5.5,
        ymin = -1, ymax = 20,
        xtick = {0,1,2,3,4},
        ytick = {0,5,10,15},
        grid = both,
        grid style = {gray!30},
        width = 12cm,
        height = 8cm,
        samples = 100,
        smooth
    ]
        \addplot[thick, domain=0:4.2] {-4*x^2 + 16*x + 2};
        
        % Points remarquables
        \addplot[only marks, mark=*, mark size=2] coordinates {(0,2) (2,18)};
        
    \end{axis}
\end{tikzpicture}
\end{center}

Déterminer avec la précision permise par le graphique :
\begin{enumerate}
\item[a)] La hauteur initiale $h_0$ à partir de laquelle la balle est lancée ;
\item[b)] La hauteur maximale $h_m$ atteinte par la balle ;
\item[c)] Le temps $t_m$ au bout duquel la hauteur maximale est atteinte ;
\item[d)] Le temps $t_1$ au bout duquel la balle touche le sol.
\end{enumerate}

\item On se propose de retrouver par le calcul les résultats de la première question.
\begin{enumerate}
\item[a)] Calculer la hauteur initiale $h_0$ à partir de laquelle la balle est lancée.
\item[b)] Calculer le temps au bout duquel la hauteur maximale $t_m$ est atteinte. En déduire la hauteur maximale $h_m$ atteinte par la balle.
\item[c)] Calculer avec la précision permise par l'aide au calcul ci-contre, le temps $t_1$ au bout duquel la balle touche le sol.

\textit{Aide au calcul :} $16^2 = 256$ ; $\sqrt{288} = 12\sqrt{2}$ et $1,5\sqrt{2} \approx 2,12$
\end{enumerate}
\end{enumerate}

\newpage

\subsection*{Exercice 2}

\subsubsection*{Partie A}

On considère une feuille de papier carrée dont l'aire est de $1024$ cm$^2$. On plie cette feuille en deux, on obtient ainsi un rectangle dont l'aire est la moitié de l'aire de départ.

\begin{center}
\begin{tikzpicture}
    % Feuille initiale
    \draw (0,0) rectangle (4,4);
    \node at (2,2) {$a_0 = 1024$ cm$^2$};
    \draw[->] (4.5,2) -- (5.5,2);
    
    % Après 1 pliage
    \draw (6,0) rectangle (10,2);
    \node at (8,1) {$a_1 = 512$ cm$^2$};
    \draw[->] (10.5,1) -- (11.5,1);
    
    % Après 2 pliages
    \draw (12,0) rectangle (14,2);
    \node at (13,1) {$a_2 = 256$ cm$^2$};
\end{tikzpicture}
\end{center}

On note $a_0$ l'aire de la feuille initiale et $a_1$ l'aire de la feuille après un pliage.

\begin{enumerate}
\item Calculer $a_1$ puis $a_2$.
\item Justifier que pour tout entier naturel $n$, $a_{n+1} = \dfrac{1}{2} a_n$.
\item En déduire l'expression de $a_n$ en fonction de $n$.
\item Calculer l'aire du carré obtenu après 10 pliages.
\end{enumerate}

\subsubsection*{Partie B -- Étude de l'aire totale}

Pour tout entier naturel $n$, on note $S_n$ la somme des $(n+1)$ aires des carrés successifs obtenus lors de $n$ pliages :
\[
S_n = a_0 + a_1 + \ldots + a_n
\]

\begin{enumerate}
\item Montrer que pour tout entier naturel $n$ :
\[
S_n = 2^{11} \left(1 - \frac{1}{2^{n+1}}\right) = 2048 \left(1 - \frac{1}{2^{n+1}}\right)
\]

\item On donne ci-contre des valeurs de termes de la suite $(S_n)$ obtenues à l'aide d'une calculatrice.
\[
\begin{array}{|c|c|}
\hline
n & S_n \\
\hline
5 & 2016 \\
10 & 2047,5 \\
15 & 2047,96875 \\
20 & 2047,9990234375 \\
\hline
\end{array}
\]
Conjecturer le comportement des termes $S_n$ lorsque $n$ tend vers $+\infty$. Interpréter ce résultat dans le contexte de l'exercice.
\end{enumerate}

\end{document}